\documentclass[aspectratio=169,10pt,english]{beamer}

\input{../../../_preambulo_beamer.tex}
\input{../../../_comandos_beamer.tex}

\selectlanguage{english}

\title{MA1001B - Statistical Modeling for Decision Making}
\subtitle{Lesson 10: Tree Diagrams and Sample-Space Decomposition}
\author{}
\institute{Tecnol\'ogico de Monterrey}
\date{}

\begin{document}

\begin{frame}[plain]
  \vspace{1.2cm}
  {\Large\bfseries MA1001B - Statistical Modeling for Decision Making\par}
  \vspace{0.7cm}
  {\large Lesson 10: Tree Diagrams and Sample-Space Decomposition\par}
  \vspace{0.7cm}
  {\color{TecRojo}\rule{\textwidth}{1.2pt}}
  \vspace{0.8cm}

  MA1001B Faculty\\[0.3cm]
  Term\\[0.3cm]
  Tecnol\'ogico de Monterrey
\end{frame}

\begin{frame}{The Sequential Probability Question}
  \begin{alertblock}{Guiding question}
    How can a two-stage experiment be decomposed into branches so that joint
    and reversed conditionals stay visible?
  \end{alertblock}

  \vspace{0.6em}
  By the end of this lesson, you can:
  \begin{itemize}
    \item draw a tree of mutually exclusive terminal paths;
    \item label branches with frequencies or probabilities;
    \item multiply along a path and add disjoint paths;
    \item explain why $P(B\mid A)$ can be a branch while $P(A\mid B)$ is not.
  \end{itemize}

  \vfill
  \scriptsize All lots and inspection flags used today are synthetic. No real company data are used.
\end{frame}

\begin{frame}{A Two-Stage Sample Space}
  \small
  \begin{center}
    \begin{tabular}{lrrr}
      \toprule
      \textbf{Lot status} & \textbf{Flagged} & \textbf{Not flagged} & \textbf{Total} \\
      \midrule
      Defective & 8 & 2 & 10 \\
      Clean & 9 & 81 & 90 \\
      \midrule
      Total & 17 & 83 & 100 \\
      \bottomrule
    \end{tabular}
  \end{center}

  \vspace{0.5em}
  \begin{exampleblock}{Synthetic inspection register}
    Stage 1 is lot status. Stage 2 is the inspection result. The four terminal
    cells are mutually exclusive paths that exhaust the 100-lot sample space.
  \end{exampleblock}
\end{frame}

\begin{frame}[fragile]{Step 1: Frequency Tree}
  \begin{columns}[T]
    \begin{column}{0.42\textwidth}
      \small
      \begin{lstlisting}
joints = {
  "def_flag": 8,
  "def_miss": 2,
  "clean_flag": 9,
  "clean_clear": 81,
}
flagged = 8 + 9
      \end{lstlisting}

      \begin{block}{Verified output}
        $100$ lots; $10$ defective\\
        Terminal counts $8$, $2$, $9$, $81$\\
        Flagged total $=17$
      \end{block}
    \end{column}
    \begin{column}{0.58\textwidth}
      \centering
      \includegraphics[width=\textwidth]{../figures/tree_diagrams_01_frequency_tree.png}
      \vspace{0.2em}
      \scriptsize Frequency tree from Step 1
    \end{column}
  \end{columns}
\end{frame}

\begin{frame}{Multiply Along Branches; Add Disjoint Paths}
  \small
  First-stage branches:
  \[
    P(D)=\frac{10}{100}=0.10,\qquad P(C)=\frac{90}{100}=0.90.
  \]
  Second-stage branches:
  \[
    P(F\mid D)=\frac{8}{10}=0.80,\qquad P(F\mid C)=\frac{9}{90}=0.10.
  \]
  One path product and the flagged total:
  \[
    P(D\cap F)=0.10\times 0.80=0.08,
  \]
  \[
    P(F)=0.08+0.09=0.17.
  \]
\end{frame}

\begin{frame}[fragile]{Step 2: Path Probabilities}
  \begin{columns}[T]
    \begin{column}{0.40\textwidth}
      \small
      \begin{lstlisting}
p_def_flag = p_d * p_f_given_d
p_flag = (
    p_def_flag
    + p_clean_flag
)
      \end{lstlisting}

      \begin{block}{Verified output}
        $P(D\cap F)=0.080000$\\
        $P(C\cap F)=0.090000$\\
        $P(F)=0.170000$\\
        Four-path sum $=1.000000$
      \end{block}
    \end{column}
    \begin{column}{0.60\textwidth}
      \centering
      \includegraphics[width=\textwidth]{../figures/tree_diagrams_02_path_probabilities.png}
      \vspace{0.2em}
      \scriptsize Terminal-path probabilities from Step 2
    \end{column}
  \end{columns}
\end{frame}

\begin{frame}{Step 3: The Reverse Conditional Is Not a Branch}
  \begin{columns}[T]
    \begin{column}{0.42\textwidth}
      \small
      $P(F\mid D)=0.800000$ is a downward branch.

      \vspace{0.5em}
      \[
        P(D\mid F)=\frac{8}{17}=0.470588.
      \]

      \vspace{0.4em}
      Seed-42 simulation of $100{,}000$ lots: $0.469551$.

      \vspace{0.5em}
      \textbf{Limit:} a time-ordered tree does not display $P(D\mid F)$ as a
      single branch. Lesson 11 names that inversion as Bayes' theorem.
    \end{column}
    \begin{column}{0.58\textwidth}
      \centering
      \includegraphics[width=\textwidth]{../figures/tree_diagrams_03_reverse_conditional.png}
      \vspace{0.2em}
      \scriptsize Branch label versus reversed conditional from Step 3
    \end{column}
  \end{columns}
\end{frame}

\begin{frame}{Concept Check}
  \begin{enumerate}
    \item Why must the four terminal counts sum to 100 in this tree?
    \item Compute $P(C\cap F)$ from the branch labels $0.90$ and $0.10$.
    \item Which probability is a single second-stage branch: $P(F\mid D)$ or
      $P(D\mid F)$?
    \item Why is $P(D\mid F)=8/17$ not equal to $P(F\mid D)=0.80$?
  \end{enumerate}

  \vspace{0.6em}
  \begin{block}{Connection to Lesson 11}
    The session can continue directly into Bayes' theorem, prior and posterior
    probabilities, and base-rate weighting.
  \end{block}

  \vfill
  \scriptsize Source: OpenStax, \textit{Introductory Business Statistics 2e},
  Chapter 3, Section 3.4. CC BY-NC-SA.
\end{frame}

\appendix



\end{document}
